%% ============================================================
%% BACK MATTER: Bibliography + Author Bios
%% ============================================================

%% ---- Bibliography ----
%% Available styles: model1-num-names (numbered), cas-model2-names (author-year)
\bibliographystyle{../model1-num-names}
%\bibliographystyle{../cas-model2-names}
%\bibliographystyle{ieeetr}

\bibliography{../cas-refs}


%% ---- Author Biographies ----

\bio{figs/Arkaprabha Laha.png}
Arkaprabha Laha received his B.Tech degree in Computer Science and Engineering from MCKV Institute of Engineering, affiliated with Maulana Abul Kalam Azad University of Technology. He is currently pursuing his Ph.D. in the Department of Computer Science and Engineering at Indian Institute of Technology (Indian School of Mines) Dhanbad. His research interests primarily include Computer Vision and related areas of Deep Learning.
\endbio

\bio{figs/Puspen Lahiri.jpg}
Puspen Lahiri received the B.Tech. degree in Computer Science and Engineering from WBUT and the M.Tech. degree in Computer Technology from Jadavpur University, where he is currently pursuing the Ph.D. degree. He is an Assistant Professor with the Department of Computer Science and Engineering, MCKV Institute of Engineering, Howrah, West Bengal, India. His research interests primarily include image restoration using deep learning, with emphasis on image processing and advanced deep learning methodologies.
\endbio

\bio{figs/Shubhomay Kundu Poddar.jpg}
Shubhomay Kundu Poddar received his Bachelor of Technology (B.Tech) degree in Computer Science and Engineering from MCKV Institute of Engineering. He has also completed a Diploma in Data Science through the BS Degree Programme in Data Science and Applications offered by Indian Institute of Technology Madras. He is currently working as a Programmer Analyst Trainee at Cognizant Technology Solutions. His research interests include Generative Artificial Intelligence, Artificial Intelligence, Machine Learning, Deep Learning, and Natural Language Processing. He is particularly interested in exploring emerging technologies and developing intelligent systems and AI-driven solutions for real-world applications.
\endbio

\bio{figs/Supriyo Ghosh.jpg}
Supriyo Ghosh received his Bachelor of Technology (B.Tech) degree in Computer Science and Engineering from MCKV Institute of Engineering in 2025. He completed a research internship at the Indian Space Research Organisation (ISRO), where he worked on computer vision and neural network applications. He is currently working as an Assistant Engineer at CodeClouds IT Solutions, specializing in backend systems engineering using Python, FastAPI, Flask, Node.js, SQL, MongoDB, Docker, Kafka, Microservices, and Redis. His work also involves integrating production-grade Machine Learning and Large Language Model (LLM) workflows. His research interests include Generative AI, Artificial Intelligence, Deep Learning, Computer Vision, Large Language Models, Edge AI, Robotics, and intelligent systems for industrial applications. He qualified the Graduate Aptitude Test in Engineering and was a finalist in the Cognizant Hackathon.
\endbio

\bio{figs/Shubhradip Pic.jpeg}
Shubradip Chowdhury received his Bachelor of Technology (B.Tech) degree in Computer Science and Engineering from MCKV Institute of Engineering. He qualified the Graduate Aptitude Test in Engineering, demonstrating strong analytical and technical skills in computer science. His research experience includes multidimensional medical image segmentation, with a focus on deep learning techniques for the accurate analysis of complex medical imaging data. His areas of interest include algorithms, artificial intelligence, computer vision, machine learning, and Java-based software development. With a strong foundation in problem-solving and programming, he aims to contribute to innovative research and advancements in intelligent computing systems through further academic and research pursuits.
\endbio

\bio{figs/H Roy.jpg}
Dr. Hiranmoy Roy is an Associate Professor in the Department of Information Technology at RCC Institute of Information Technology, Kolkata, India, and earned his Ph.D. in Computer Science and Engineering from Jadavpur University under Prof. Debotosh Bhattacharjee, focusing on feature extraction for heterogeneous face recognition. With over 23 years of academic and research experience, his interests include Image Processing, Computer Vision, Machine Learning, Deep Learning, Artificial Intelligence, and Quantum Computing. He has published journal papers in leading journals such as IEEE TIFS, IEEE TPAMI, and Elsevier Information Sciences. Dr. Roy currently works on interpretable AI models, human visual system inspired networks, quantum domain image descriptors, and AI-enabled solutions for healthcare, agriculture, and digital forensics.
\endbio

\bio{figs/D Bhattacharjee.jpg}
Debotosh Bhattacharjee, a full professor in the Department of Computer Science and Engineering at Jadavpur University, has research interests includes applications of machine learning techniques in Biometrics and Diagnostic Image Analysis. He has authored or co-authored more than 400 journal articles, conference proceedings publications, and book chapters.  Two US, one Australian, two German, and one Indian patents have been granted on his works. Prof. Bhattacharjee has been awarded sponsored projects by the Government of India funding agencies, totaling approximately INR 30 million. He has served as a visiting professor and postdoctoral researcher at various universities abroad, including those in the Netherlands, Portugal, Russia, Slovenia, the Czech Republic, the UK, and Germany. Prof. Bhattacharjee is actively involved in Outcome-Based Education (OBE) and accreditation, serving as a resource person and empaneled evaluator for NBA, NAAC, and AICTE. He is on the editorial board of several journals. He has organized nine international conferences as a general chair with overall responsibility. He is a Fellow of the WBAST, IETE, a Life Member of ISTE, and a Senior Member of IEEE (USA).
\endbio
